% Options for packages loaded elsewhere
% Options for packages loaded elsewhere
\PassOptionsToPackage{unicode}{hyperref}
\PassOptionsToPackage{hyphens}{url}
\PassOptionsToPackage{dvipsnames,svgnames,x11names}{xcolor}
%
\documentclass[
  spanish,
  a4paper,
  oneside]{scrbook}
\usepackage{xcolor}
\usepackage[top=25mm,bottom=25mm,left=30mm,right=20mm]{geometry}
\usepackage{amsmath,amssymb}
\setcounter{secnumdepth}{5}
\usepackage{iftex}
\ifPDFTeX
  \usepackage[T1]{fontenc}
  \usepackage[utf8]{inputenc}
  \usepackage{textcomp} % provide euro and other symbols
\else % if luatex or xetex
  \usepackage{unicode-math} % this also loads fontspec
  \defaultfontfeatures{Scale=MatchLowercase}
  \defaultfontfeatures[\rmfamily]{Ligatures=TeX,Scale=1}
\fi
\usepackage{lmodern}
\ifPDFTeX\else
  % xetex/luatex font selection
\fi
% Use upquote if available, for straight quotes in verbatim environments
\IfFileExists{upquote.sty}{\usepackage{upquote}}{}
\IfFileExists{microtype.sty}{% use microtype if available
  \usepackage[]{microtype}
  \UseMicrotypeSet[protrusion]{basicmath} % disable protrusion for tt fonts
}{}
\makeatletter
\@ifundefined{KOMAClassName}{% if non-KOMA class
  \IfFileExists{parskip.sty}{%
    \usepackage{parskip}
  }{% else
    \setlength{\parindent}{0pt}
    \setlength{\parskip}{6pt plus 2pt minus 1pt}}
}{% if KOMA class
  \KOMAoptions{parskip=half}}
\makeatother
% Make \paragraph and \subparagraph free-standing
\makeatletter
\ifx\paragraph\undefined\else
  \let\oldparagraph\paragraph
  \renewcommand{\paragraph}{
    \@ifstar
      \xxxParagraphStar
      \xxxParagraphNoStar
  }
  \newcommand{\xxxParagraphStar}[1]{\oldparagraph*{#1}\mbox{}}
  \newcommand{\xxxParagraphNoStar}[1]{\oldparagraph{#1}\mbox{}}
\fi
\ifx\subparagraph\undefined\else
  \let\oldsubparagraph\subparagraph
  \renewcommand{\subparagraph}{
    \@ifstar
      \xxxSubParagraphStar
      \xxxSubParagraphNoStar
  }
  \newcommand{\xxxSubParagraphStar}[1]{\oldsubparagraph*{#1}\mbox{}}
  \newcommand{\xxxSubParagraphNoStar}[1]{\oldsubparagraph{#1}\mbox{}}
\fi
\makeatother

\usepackage{color}
\usepackage{fancyvrb}
\newcommand{\VerbBar}{|}
\newcommand{\VERB}{\Verb[commandchars=\\\{\}]}
\DefineVerbatimEnvironment{Highlighting}{Verbatim}{commandchars=\\\{\}}
% Add ',fontsize=\small' for more characters per line
\usepackage{framed}
\definecolor{shadecolor}{RGB}{241,243,245}
\newenvironment{Shaded}{\begin{snugshade}}{\end{snugshade}}
\newcommand{\AlertTok}[1]{\textcolor[rgb]{0.68,0.00,0.00}{#1}}
\newcommand{\AnnotationTok}[1]{\textcolor[rgb]{0.37,0.37,0.37}{#1}}
\newcommand{\AttributeTok}[1]{\textcolor[rgb]{0.40,0.45,0.13}{#1}}
\newcommand{\BaseNTok}[1]{\textcolor[rgb]{0.68,0.00,0.00}{#1}}
\newcommand{\BuiltInTok}[1]{\textcolor[rgb]{0.00,0.23,0.31}{#1}}
\newcommand{\CharTok}[1]{\textcolor[rgb]{0.13,0.47,0.30}{#1}}
\newcommand{\CommentTok}[1]{\textcolor[rgb]{0.37,0.37,0.37}{#1}}
\newcommand{\CommentVarTok}[1]{\textcolor[rgb]{0.37,0.37,0.37}{\textit{#1}}}
\newcommand{\ConstantTok}[1]{\textcolor[rgb]{0.56,0.35,0.01}{#1}}
\newcommand{\ControlFlowTok}[1]{\textcolor[rgb]{0.00,0.23,0.31}{\textbf{#1}}}
\newcommand{\DataTypeTok}[1]{\textcolor[rgb]{0.68,0.00,0.00}{#1}}
\newcommand{\DecValTok}[1]{\textcolor[rgb]{0.68,0.00,0.00}{#1}}
\newcommand{\DocumentationTok}[1]{\textcolor[rgb]{0.37,0.37,0.37}{\textit{#1}}}
\newcommand{\ErrorTok}[1]{\textcolor[rgb]{0.68,0.00,0.00}{#1}}
\newcommand{\ExtensionTok}[1]{\textcolor[rgb]{0.00,0.23,0.31}{#1}}
\newcommand{\FloatTok}[1]{\textcolor[rgb]{0.68,0.00,0.00}{#1}}
\newcommand{\FunctionTok}[1]{\textcolor[rgb]{0.28,0.35,0.67}{#1}}
\newcommand{\ImportTok}[1]{\textcolor[rgb]{0.00,0.46,0.62}{#1}}
\newcommand{\InformationTok}[1]{\textcolor[rgb]{0.37,0.37,0.37}{#1}}
\newcommand{\KeywordTok}[1]{\textcolor[rgb]{0.00,0.23,0.31}{\textbf{#1}}}
\newcommand{\NormalTok}[1]{\textcolor[rgb]{0.00,0.23,0.31}{#1}}
\newcommand{\OperatorTok}[1]{\textcolor[rgb]{0.37,0.37,0.37}{#1}}
\newcommand{\OtherTok}[1]{\textcolor[rgb]{0.00,0.23,0.31}{#1}}
\newcommand{\PreprocessorTok}[1]{\textcolor[rgb]{0.68,0.00,0.00}{#1}}
\newcommand{\RegionMarkerTok}[1]{\textcolor[rgb]{0.00,0.23,0.31}{#1}}
\newcommand{\SpecialCharTok}[1]{\textcolor[rgb]{0.37,0.37,0.37}{#1}}
\newcommand{\SpecialStringTok}[1]{\textcolor[rgb]{0.13,0.47,0.30}{#1}}
\newcommand{\StringTok}[1]{\textcolor[rgb]{0.13,0.47,0.30}{#1}}
\newcommand{\VariableTok}[1]{\textcolor[rgb]{0.07,0.07,0.07}{#1}}
\newcommand{\VerbatimStringTok}[1]{\textcolor[rgb]{0.13,0.47,0.30}{#1}}
\newcommand{\WarningTok}[1]{\textcolor[rgb]{0.37,0.37,0.37}{\textit{#1}}}

\usepackage{longtable,booktabs,array}
\newcounter{none} % for unnumbered tables
\usepackage{calc} % for calculating minipage widths
% Correct order of tables after \paragraph or \subparagraph
\usepackage{etoolbox}
\makeatletter
\patchcmd\longtable{\par}{\if@noskipsec\mbox{}\fi\par}{}{}
\makeatother
% Allow footnotes in longtable head/foot
\IfFileExists{footnotehyper.sty}{\usepackage{footnotehyper}}{\usepackage{footnote}}
\makesavenoteenv{longtable}
\usepackage{graphicx}
\makeatletter
\newsavebox\pandoc@box
\newcommand*\pandocbounded[1]{% scales image to fit in text height/width
  \sbox\pandoc@box{#1}%
  \Gscale@div\@tempa{\textheight}{\dimexpr\ht\pandoc@box+\dp\pandoc@box\relax}%
  \Gscale@div\@tempb{\linewidth}{\wd\pandoc@box}%
  \ifdim\@tempb\p@<\@tempa\p@\let\@tempa\@tempb\fi% select the smaller of both
  \ifdim\@tempa\p@<\p@\scalebox{\@tempa}{\usebox\pandoc@box}%
  \else\usebox{\pandoc@box}%
  \fi%
}
% Set default figure placement to htbp
\def\fps@figure{htbp}
\makeatother


% definitions for citeproc citations
\NewDocumentCommand\citeproctext{}{}
\NewDocumentCommand\citeproc{mm}{%
  \begingroup\def\citeproctext{#2}\cite{#1}\endgroup}
\makeatletter
 % allow citations to break across lines
 \let\@cite@ofmt\@firstofone
 % avoid brackets around text for \cite:
 \def\@biblabel#1{}
 \def\@cite#1#2{{#1\if@tempswa , #2\fi}}
\makeatother
\newlength{\cslhangindent}
\setlength{\cslhangindent}{1.5em}
\newlength{\csllabelwidth}
\setlength{\csllabelwidth}{3em}
\newenvironment{CSLReferences}[2] % #1 hanging-indent, #2 entry-spacing
 {\begin{list}{}{%
  \setlength{\itemindent}{0pt}
  \setlength{\leftmargin}{0pt}
  \setlength{\parsep}{0pt}
  % turn on hanging indent if param 1 is 1
  \ifodd #1
   \setlength{\leftmargin}{\cslhangindent}
   \setlength{\itemindent}{-1\cslhangindent}
  \fi
  % set entry spacing
  \setlength{\itemsep}{#2\baselineskip}}}
 {\end{list}}
\usepackage{calc}
\newcommand{\CSLBlock}[1]{\hfill\break\parbox[t]{\linewidth}{\strut\ignorespaces#1\strut}}
\newcommand{\CSLLeftMargin}[1]{\parbox[t]{\csllabelwidth}{\strut#1\strut}}
\newcommand{\CSLRightInline}[1]{\parbox[t]{\linewidth - \csllabelwidth}{\strut#1\strut}}
\newcommand{\CSLIndent}[1]{\hspace{\cslhangindent}#1}

\ifLuaTeX
\usepackage[bidi=basic,shorthands=off]{babel}
\else
\usepackage[bidi=default,shorthands=off]{babel}
\fi
\ifLuaTeX
  \usepackage{selnolig} % disable illegal ligatures
\fi


\setlength{\emergencystretch}{3em} % prevent overfull lines

\providecommand{\tightlist}{%
  \setlength{\itemsep}{0pt}\setlength{\parskip}{0pt}}



 


\AtBeginDocument{\hypersetup{linkcolor=black}}
% ---------- Paquetes tipograficos y microajustes ----------
\usepackage{microtype}        % mejor interletrado y justificado
\usepackage{csquotes}         % comillas tipograficas
\usepackage{iftex}
\ifPDFTeX\else
  % Seleccion de fuentes solo cuando XeTeX/LuaTeX esta disponible.
  \IfFontExistsTF{TeX Gyre Termes}{
    \setmainfont{TeX Gyre Termes}
  }{
    \IfFontExistsTF{Latin Modern Roman}{\setmainfont{Latin Modern Roman}}{}
  }
  \IfFontExistsTF{TeX Gyre Heros}{
    \setsansfont{TeX Gyre Heros}
  }{
    \IfFontExistsTF{Latin Modern Sans}{\setsansfont{Latin Modern Sans}}{}
  }
  \IfFontExistsTF{TeX Gyre Cursor}{
    \setmonofont{TeX Gyre Cursor}
  }{
    \IfFontExistsTF{Latin Modern Mono}{\setmonofont{Latin Modern Mono}}{}
  }
\fi
\usepackage{enumitem}         % listas compactas
\setlist{itemsep=.2em, topsep=.2em}

% ---------- KOMA-Script: estilo de titulos y espaciado ----------
\KOMAoptions{
  headings=big,
  parskip=half,
  fontsize=12pt,
  appendixprefix=true
}
\setlength{\parindent}{1.5em}
\setlength{\parskip}{0.6em}
\usepackage{indentfirst}
\usepackage{setspace}         % Control del interlineado del documento
\setstretch{1.5}

% ---------- Encabezados y pies (scrlayer-scrpage) ----------
\usepackage[automark,headsepline]{scrlayer-scrpage}
\clearpairofpagestyles
\automark[chapter]{chapter}
\ihead{\pagemark}
\ohead{\itshape\headmark}
\setheadsepline{0.4pt}
\renewcommand*{\chaptermarkformat}{}
\renewcommand*{\chapterpagestyle}{scrheadings}
\pagestyle{scrheadings}
\makeatletter
\let\ps@plain\ps@scrheadings
\makeatother
\cfoot{}

% ---------- Hipervinculos mas sobrios ----------
\usepackage{hyperref}
\hypersetup{
  colorlinks=true,
  linkcolor=black,
  citecolor=[rgb]{0.05,0.2,0.5},
  urlcolor=[rgb]{0.05,0.2,0.5},
  pdfauthor={\@author},
  pdftitle={\@title}
}

% ---------- Leyendas de figuras/tablas ----------
\usepackage[labelfont=bf,textfont=it]{caption}
\captionsetup{
  skip=8pt
}

% ---------- Tabla de contenidos ----------
\KOMAoptions{toc=graduated}
\RedeclareSectionCommand[tocnumwidth=3em]{chapter}
\RedeclareSectionCommand[tocindent=3.25em,tocnumwidth=2.8em]{section}
\RedeclareSectionCommand[tocindent=6.5em,tocnumwidth=2.5em]{subsection}
\makeatletter
\newcommand*{\tocdotfill}{\leavevmode\leaders\hbox to .6em{\hss.\hss}\hfill}
\makeatother
\RedeclareSectionCommand[toclinefill=\tocdotfill]{chapter}
\RedeclareSectionCommand[toclinefill=\tocdotfill]{section}
\RedeclareSectionCommand[toclinefill=\tocdotfill]{subsection}
\renewcommand*{\contentsname}{Tabla de contenido}
\setkomafont{chapterentry}{\normalfont}
\setkomafont{chapterentrypagenumber}{\normalfont}

% ---------- Viudas/Huerfanas y cortes de pagina ----------
\clubpenalty=10000
\widowpenalty=10000
\displaywidowpenalty=10000

% ---------- Entorno abstract para scrbook ----------
\providecommand{\abstractname}{Resumen}
\makeatletter
\@ifundefined{abstract}{
  \newenvironment{abstract}{
    \cleardoublepage
    \thispagestyle{plain}
    \null\vfill
    \begin{center}
      {\bfseries\Large \abstractname\par}
    \end{center}\vspace{1em}
    \begingroup
  }{
    \par\endgroup
    \vfill\null
    \cleardoublepage
  }
}{}
\makeatother

% ---------- Soporte de subtitulo desde YAML ----------
\makeatletter
\providecommand{\subtitle}[1]{\gdef\@subtitle{#1}}
\providecommand{\@subtitle}{}
\makeatother

\makeatletter
\newcommand{\facso@iflist}[4]{%
  \begingroup
    \def\addvspace##1{}%
    \def\numberline##1{##1}%
    \def\facso@target{#2}%
    \newcount\facso@entrycount
    \facso@entrycount=0\relax
    \long\def\contentsline##1##2##3{%
      \def\facso@this{##1}%
      \ifx\facso@this\facso@target
        \advance\facso@entrycount\@ne
      \fi
    }%
    \InputIfFileExists{\jobname.#1}{}{}%
  \endgroup
  \ifnum\facso@entrycount>0\relax
    #3%
  \else
    #4%
  \fi}
\renewcommand{\listoffigures}{%
  \facso@iflist{lof}{figure}{%
    \cleardoublepage
    \phantomsection
    \addchap*{Listado de Figuras}%
    \thispagestyle{scrheadings}%
    \markboth{Listado de Figuras}{Listado de Figuras}%
    \addcontentsline{toc}{chapter}{Listado de Figuras}%
    \@starttoc{lof}%
    \cleardoublepage
    \automark[chapter]{chapter}%
  }{}}
\renewcommand{\listoftables}{%
  \facso@iflist{lot}{table}{%
    \cleardoublepage
    \phantomsection
    \addchap*{Listado de Tablas}%
    \thispagestyle{scrheadings}%
    \markboth{Listado de Tablas}{Listado de Tablas}%
    \addcontentsline{toc}{chapter}{Listado de Tablas}%
    \@starttoc{lot}%
    \cleardoublepage
    \automark[chapter]{chapter}%
  }{}}
\makeatother

% --- Desactivar portada y abstract automaticos de Pandoc (PDF) ---
\AtBeginDocument{\let\maketitle\relax}
\renewenvironment{abstract}{}{}
\providecommand{\appendixname}{}
\providecommand{\appendixtocname}{}
\providecommand{\appendixpagename}{}
\renewcommand*{\appendixname}{Anexo}
\renewcommand*{\appendixtocname}{Anexos}
\renewcommand*{\appendixpagename}{Anexos}
\makeatletter
\@ifpackageloaded{bookmark}{}{\usepackage{bookmark}}
\makeatother
\makeatletter
\@ifpackageloaded{caption}{}{\usepackage{caption}}
\AtBeginDocument{%
\ifdefined\contentsname
  \renewcommand*\contentsname{Tabla de contenidos}
\else
  \newcommand\contentsname{Tabla de contenidos}
\fi
\ifdefined\listfigurename
  \renewcommand*\listfigurename{Listado de Figuras}
\else
  \newcommand\listfigurename{Listado de Figuras}
\fi
\ifdefined\listtablename
  \renewcommand*\listtablename{Listado de Tablas}
\else
  \newcommand\listtablename{Listado de Tablas}
\fi
\ifdefined\figurename
  \renewcommand*\figurename{Lista de figuras}
\else
  \newcommand\figurename{Lista de figuras}
\fi
\ifdefined\tablename
  \renewcommand*\tablename{Lista de tablas}
\else
  \newcommand\tablename{Lista de tablas}
\fi
}
\@ifpackageloaded{float}{}{\usepackage{float}}
\floatstyle{ruled}
\@ifundefined{c@chapter}{\newfloat{codelisting}{h}{lop}}{\newfloat{codelisting}{h}{lop}[chapter]}
\floatname{codelisting}{Listado}
\newcommand*\listoflistings{\listof{codelisting}{Listado de Listados}}
\makeatother
\makeatletter
\makeatother
\makeatletter
\@ifpackageloaded{caption}{}{\usepackage{caption}}
\@ifpackageloaded{subcaption}{}{\usepackage{subcaption}}
\makeatother
\usepackage{bookmark}
\IfFileExists{xurl.sty}{\usepackage{xurl}}{} % add URL line breaks if available
\urlstyle{same}
\hypersetup{
  pdftitle={Reportes FACSO - Plantilla Quarto},
  pdfauthor={Nombre Apellido},
  pdflang={es},
  colorlinks=true,
  linkcolor={Maroon},
  filecolor={Maroon},
  citecolor={Blue},
  urlcolor={Blue},
  pdfcreator={LaTeX via pandoc}}


\title{Reportes FACSO - Plantilla Quarto}
\usepackage{etoolbox}
\makeatletter
\providecommand{\subtitle}[1]{% add subtitle to \maketitle
  \apptocmd{\@title}{\par {\large #1 \par}}{}{}
}
\makeatother
\subtitle{Tesis, informes e investigaciones}
\author{Nombre Apellido}
\date{11 de septiembre de 2026}
\begin{document}
\frontmatter
\maketitle

% --- title-pdf.tex personalizado para portada formal ---
\makeatletter
\providecommand{\subtitle}[1]{\gdef\@subtitle{#1}}
\providecommand{\@subtitle}{}
\providecommand{\frontmattercontext}{}
\providecommand{\advisorname}{}
\providecommand{\advisorlabel}{Profesor guia:}
\providecommand{\frontmatterlocation}{}
\InputIfFileExists{includes/cover-config.tex}{}{}

\newcommand{\PrintTitle}{%
  {\sffamily\bfseries\fontsize{24pt}{28pt}\selectfont \@title\par}%
}
\newcommand{\PrintSubtitle}{%
  \begingroup
  \edef\temp{\detokenize{\@subtitle}}%
  \ifx\temp\empty\relax
    % sin subtitulo
  \else
    {\sffamily\bfseries\large \@subtitle\par}%
  \fi
  \endgroup
}
\newcommand{\PrintAuthor}{%
  {\Large\bfseries \@author\par}%
}
\newcommand{\PrintDate}{%
  {\small \@date\par}%
}
\makeatother

\begin{titlepage}
\thispagestyle{empty}
\begin{center}
\vspace*{10mm}

% Logo o imagen institucional

\includegraphics[width=0.35\textwidth]{assets/cover.png}\par
\vspace{12mm}

% Titulo y subtitulo
\PrintTitle
\vspace{6mm}
\PrintSubtitle

\vspace{22mm}
\begingroup
\edef\temp{\detokenize{\frontmattercontext}}%
\ifx\temp\empty\relax
  % sin contexto adicional
\else
  {\normalsize \frontmattercontext\par}
\fi
\endgroup

\vspace{18mm}
\PrintAuthor

\vspace{16mm}
\begingroup
\edef\temp{\detokenize{\advisorname}}%
\ifx\temp\empty\relax
  % sin tutor
\else
  \rule{0.45\textwidth}{0.4pt}\par
  {\small \advisorlabel\ \advisorname\par}
\fi
\endgroup

\vfill
\begingroup
\edef\temp{\detokenize{\frontmatterlocation}}%
\ifx\temp\empty\relax
  % sin ubicacion
\else
  {\small \frontmatterlocation\par}
\fi
\endgroup
\PrintDate

\end{center}
\end{titlepage}

% Preliminares (numeros romanos) y estilo simple
\frontmatter
\pagestyle{scrheadings}

\setcounter{tocdepth}{2} % (o 1/3 segun prefieras)
\tableofcontents
\listoffigures
\listoftables


\mainmatter
\bookmarksetup{startatroot}

\chapter*{Abstract}\label{abstract}
\addcontentsline{toc}{chapter}{Abstract}

\markboth{Abstract}{Abstract}

\small\textbf{Palabras clave:} Sociologia, educacion, Chile \normalsize

Resumen en español (200-300). Expón propósito, método, resultados clave
y conclusiones. Puede incluir citas como
(\citeproc{ref-xie2017bookdown}{\textbf{xie2017bookdown?}}) y se
actualizarán automáticamente.

\bookmarksetup{startatroot}

\chapter*{Prefacio}\label{prefacio}
\addcontentsline{toc}{chapter}{Prefacio}

\markboth{Prefacio}{Prefacio}

\mainmatter

\bookmarksetup{startatroot}

\chapter{Introducción}\label{introducciuxf3n}

\subsection{Introduccion}\label{introduccion}

La preocupación por la preparación ciudadana de los jóvenes en el
cotexto local ha impulsado diversas iniciativas curriculares y
legislativas durannte el ultimo tiempo. Sin embargo, la evidencia ha
dado cuennta de un escenario con importantes desafíos. De acuerdo con el
Estudio Internacional de Educación Cívica y Formación Ciudadana
(ICCSref2016), los estudiantes chilenos de 8° básico alcanzaron un
puntaje promedio de 482 puntos en conocimiento cívico, situándose
significativamente por debajo de la media internacional de 517 puntos.
Este rendimiento se mantuvo estancado entre los ciclos 2009 y 2016, a
diferencia de otros países de la región como Colombia y México que
registraron avances significativos. En cuanto al compromiso cívico,
existe una brecha marcada entre el entorno escolar y el espacio público
exterior. Si bien dentro de las escuelas se registran altos y buenos
niveles de participación; como el voto en elecciones escolares y el
involucramiento en la toma de decisiones, fuera del colegio la discusión
sobre temas políticos contingentes y la participación en redes sociales
caen por debajo del promedio internacional. Del mismo modo, las
expectativas de participación electoral futura en la adultez se
mantienen congeladas en 50 puntos (bajo la media internacional de 51)
desde 2009.

Si bien la Ley N° 20.911, promulgada el año 2016 mandató la
implementación de Planes de Formación Ciudadana en todos los colegios
teniedo altas expectativas sobre el rol de la escuela en la promocion de
la ciudadania activa, la evidencia empírica invita a matizar el alcance
real de la institución escolar. El análisis de
(\citeproc{ref-pdf_d}{\textbf{pdf\_d?}}) sobre los datos de ICCS revela
que la escuela incide de manera diferenciada en la formación ciudadana.
Gran parte de la diferencia de los puntajes en conocimento civico se
debe a diferencias entre escuelas; pero estas mismas diferencias on
explican de manera significativa en las actitudes y disposiciones de los
estudiantes.

Esto evidencia que la escuela es efectiva transmitiendo contenidos
cognitivos, pero tiene capacidades limitadas para moldear valores y
actitudes, los cuales se encuentran fuertemente condicionados por la
familia, los pares y el entorno comunitario. Sin embargo, la escuela
chilena es carazteristica por su particular economía moral,
caracterizada por su fuerte orientacion meritocratica, inculcando la
idea de que el rendimiento y el éxito académico responden exclusivamente
al talento y esfuerzo individual (Canales, 2022; Crawford 2010;
Wiederkehr et al.~2015).

La promocion de los valores meritocratticos en la escuela chilena ha
sido señalada como un factor que erosiona la noción del bien común y la
solidaridad, promoviendo la competencia individual y la
auto-responsabilización frente al desempeño escolar (Batruch et
al.~2023; Castillo 2024; Darnon et al.~2018; Wiederkehr et al.~2015). La
internalización de estas creencias meritocráticas puede llevar a los
estudiantes a justificar las desigualdades sociales y a disminuir su
disposición hacia la acción colectiva como medio de transformación
social.

De este modo, el presente trabajo pretende analizar la relacion
contradicción de fondo; los esfuerzos de la escuela por formar
ciudadanías democráticas compiten cotidianamente con una socialización
implícita orientada a la competencia y al mérito individual que erosiona
la noción del bien común. Analizar la adscripción al mérito resulta, por
lo tanto, indispensable para explicar cómo se configuran realmente las
actitudes y comportamientos cívicos de las nuevas generaciones.

\subsection{Antecedentes}\label{antecedentes}

Chile figura como un escenario sociopolítico y económico caracterizado
por desigualdades persistentes, consolidadas a partir del despliegue de
las reformas neoliberales implementadas desde los años ochenta
(\citeproc{ref-palet_desiguales_2017}{Palet et~al., 2017}). La
literatura ha descrito consistentemente el caso de chile como un régimen
de bienestar ``liberal-informal'' debido a la particular
mercantilización de las esferas de reproducción
(\citeproc{ref-ferre_welfare_2023}{Ferre, 2023}). El país exhibe altos
niveles de desigualdad económica que se ubican de manera sistemática
entre los más altos tanto de latinoamerica como de la OCDE
(\citeproc{ref-palet_desiguales_2017}{Palet et~al., 2017}). Pese a esta
persistencia de la desigualdad y al hecho de qu , la subjetividad moral
de la población es caracteristica por una fuerte y paradójica adhesión
cultural a la ideología del mérito personal
(\citeproc{ref-castillo_meritocracia_2019}{Castillo et~al., 2019}). Las
conclusiones obtenindas a partir de los aalisis en las encuestas panel
como el Estudio Social Longitudinal de Chile (ELSOC) reportan de manera
consistente que cerca del 75\% al 80\% de los chilenos considera que
``el trabajo duro es un medio esencial para surgir en la vida''
(\citeproc{ref-castillo_meritocracia_2019}{Castillo et~al., 2019}). Sin
embargo, esta alta valoración del ideal meritocrático coexiste con una
con una contradicción a nivel subjetivo; solo una porcion menor de la
población (aproximadamente un 26\% a 30\%) siente que en la práctica las
personas en Chile son efectivamente recompensadas por sus esfuerzos
(\citeproc{ref-castillo_meritocracia_2019}{Castillo et~al., 2019};
\citeproc{ref-luismaldonado_meritocracia_2021}{Maldonado et~al., 2021}).

Esta paradoja meritocrática convive con una particular configuracion en
el campo relacionado a la ciudadania y la socialización política de las
generacionnes mmas jovenes. El sistema educativo chileno ha operado
históricamente bajo las reglas de un cuasi-mercado escolar altamente
segregado por nivel socioeconómico, estructurado a partir del
financiamiento compartido, el incentivo a la selección de alumnos y la
competencia de matrículas basada en rendición de cuentas
(\citeproc{ref-verger_recontextualizacion_2016}{Verger et~al., 2016}).
En este contexto institucional, la escuela chilena ha promovido
intensamente lo que desde la sociología de la educacion se denonimna
como el ``mito meritocrático''
(\citeproc{ref-canalesceron_sujeto_2022}{Canales Cerón et~al., 2022};
\citeproc{ref-crawford_schooling_2010}{Crawford, 2010}), infundiendo en
los estudiantes la creencia de que el rendimiento académico es un
reflejo transparente y neutro de sus propios talentos y esfuerzos
individuales (\citeproc{ref-wiederkehr_belief_2015}{Wiederkehr et~al.,
2015}). De forma simultánea, Chile enfrenta un prolongado déficit en
participación y competencias en cuanto a formacion cívica y ciudadanía,
ambitos en los que posee una de las tasas de votación juvenil más bajas
a nivel mundial tras la implementación del voto voluntario
(\citeproc{ref-somma_no_2020}{Somma et~al., 2020}). Sumado a lo
anterior, los estudiantes de octavo básico en pruebas internacionales
como el Estudio Internacional de Educación Cívica y Ciudadana (ICCS
2016) promediaron un puntaje de 482 puntos, ubicándose
significativamente por debajo del estándar global de 517 puntos
(\citeproc{ref-schulz_becoming_2018}{Schulz et~al., 2018}).

A partir de la reincorporación obligatoria en 2020 de la asignatura de
Formación Ciudadana en la enseñanza media (Ley N° 20.911), una de las
principales interrogantes que surgen es por la realizibilidad de poder
realmente formar ciudadaníos comprometidos con los principios de
equidad, solidaridad y cohesión, en un entorno educativo cuya economía
moral sigue estando anclada en la competencia.

La literatura ha anteriormente señalado una profunda tensión estructural
entre los principios de la meritocracia escolar y los objetivos de la
formación ciudadana (\citeproc{ref-castillo_socialization_2024}{Castillo
et~al., 2024}). Mientras que la ideología meritocrática define la
distribución de bienes y recompensas en función del talento, la
competencia y el esfuerzo individual
(\citeproc{ref-castillo_socialization_2024}{Castillo et~al., 2024};
\citeproc{ref-young_triunfo_}{Young, s.~f.}), la educación ciudadana
busca cultivar la pertenencia cívica, la solidaridad, el bien común y la
cooperación comunitaria. Cuando la institución escolar sobrevalora el
logro individual a través del denominado ``mito o cuento
meritocrático'', genera en los estudiantes de menor estatus
socioeconómico dinámicas de auto-responsabilización frente al desempeño
o fracaso escolar (\citeproc{ref-darnon_where_2018}{Darnon et~al.,
2018}; \citeproc{ref-wiederkehr_belief_2015}{Wiederkehr et~al., 2015}).

Esta atribución interna conduce a que los alumnos desaventajados
asimilen su posición como merecida, invisibilizando las barreras
socioestructurales de origen, promoviendo sentimientos de culpa y
autodesvalorización (\citeproc{ref-wiederkehr_belief_2015}{Wiederkehr
et~al., 2015}), y debilitando la disposición hacia la acción colectiva
como vía legítima de transformación social
(\citeproc{ref-batruch_belief_2023}{Batruch et~al., 2023};
\citeproc{ref-palet_desiguales_2017}{Palet et~al., 2017}). Asimismo, las
creencias meritocráticas operan como un canal directo de legitimación de
las desigualdades que altera las preferencias normativas de los jóvenes
(\citeproc{ref-batruch_belief_2023}{Batruch et~al., 2023};
\citeproc{ref-castillo_socialization_2024}{Castillo et~al., 2024}).

La evidencia empírica en población escolar chilena demuestra que
percibir que la meritocracia funciona en la sociedad o en la escuela
incrementa la adhesión a las preferencias de justicia de mercado, es
decir, la aceptación moral de que el acceso a servicios sociales
esenciales como salud, educación y pensiones deba estar condicionado a
la capacidad de pago individual
(\citeproc{ref-castillo_socialization_2024}{Castillo et~al., 2024}). Del
mismo modo, la Creencia en la Meritocracia Escolar (BSM) actúa como una
ideología de justificación del sistema que reduce la percepción de
injusticia frente a la estratificación social y disminuye el apoyo a
políticas redistributivas o de acción afirmativa
(\citeproc{ref-batruch_belief_2023}{Batruch et~al., 2023};
\citeproc{ref-darnon_where_2018}{Darnon et~al., 2018};
\citeproc{ref-wiederkehr_belief_2015}{Wiederkehr et~al., 2015}). De este
modo, la internalización temprana del mérito sustituye la concepción de
los servicios básicos como derechos ciudadanos igualitarios por una
aceptación de la provisión diferenciada de bienestar según el éxito
mercantil (\citeproc{ref-castillo_socialization_2024}{Castillo et~al.,
2024}).

En contraste con la sobrevaloración del logro individualista, la
percepción de la escuela como un espacio procedimental y
distributivamente justo ejerce un impacto positivo en la conducta cívica
de los estudiantes (\citeproc{ref-resh_sense_2010}{Resh, 2010};
\citeproc{ref-resh_sense_2017}{Resh \& Sabbagh, 2017}). La investigación
demuestra que cuando los alumnos evalúan que las notas y las relaciones
interpersonales con sus profesores responden a criterios de justicia y
reconocimiento, disminuyen los comportamientos de indisciplina y
violencia escolar (\citeproc{ref-resh_sense_2017}{Resh \& Sabbagh,
2017}), e incrementan significativamente su participación en proyectos
prosociales, la vinculación comunitaria y el voluntariado
(\citeproc{ref-resh_sense_2017}{Resh \& Sabbagh, 2017})

El origen socioeconómico de los estudiantes, y el acceso al capital
cultural en el hogar constituye uno de los predictores más consistentes
del conocimiento cívico
(\citeproc{ref-abendschon_political_2017}{Abendschön \& Tausendpfund,
2017}; \citeproc{ref-_pdf_c}{{«({PDF}) {Desigualdad} y {Conocimiento
Cívico}»}, s.~f.}; \citeproc{ref-schulz_becoming_2018}{Schulz et~al.,
2018}). Los estudiantes provenientes de familias de mayor estatus
obtienen consistentemente mayores puntajes en las evaluaciones de
conocimiento cívico; al mismo tiempo que manifiestan mayores
expectativas de participación política y cívica en la adultez
(\citeproc{ref-castillo_social_2014}{J. Castillo et~al., 2014};
\citeproc{ref-jerome_systematic_2024}{Jerome et~al., 2024};
\citeproc{ref-_pdf_c}{{«({PDF}) {Desigualdad} y {Conocimiento Cívico}»},
s.~f.}). La frecuencia con la que se discuten temas sociales en el hogar
y el interés político de los padres son factores de peso para predecir
el comportamiento cívico y el aprendizaje de las nuevas generaciones
(\citeproc{ref-flanagan_adolescents_2014}{Flanagan et~al., 2014};
\citeproc{ref-isac_contribution_2013}{{Isac et~al.}, 2013};
\citeproc{ref-jerome_systematic_2024}{Jerome et~al., 2024};
\citeproc{ref-lee_processes_2012}{Lee et~al., 2012}).

A nivel escolar, el clima expuesto a la discusión representa una de las
herramientas didacticas de mayor impacto para la formacion ciudadana
(\citeproc{ref-campbell_voice_2008}{Campbell, 2008};
\citeproc{ref-isac_contribution_2013}{{Isac et~al.}, 2013};
\citeproc{ref-jerome_systematic_2024}{Jerome et~al., 2024};
\citeproc{ref-schulz_becoming_2018}{Schulz et~al., 2018}). Aquellas
aulas de clase donde los profesores incentivan el diálogo
critico/reflexivo respecto a múltiples perspectivas ante temas
controvertidos y respetan la expresión de opiniones diversas potencian
sustantivamente el conocimiento cívico de los estudiantes
(\citeproc{ref-jerome_systematic_2024}{Jerome et~al., 2024};
\citeproc{ref-schulz_becoming_2018}{Schulz et~al., 2018}). Además, este
ambiente deliberativo fortalece la autoeficacia cívica de los jóvenes,
la conciencia crítica, sus actitudes hacia la igualdad de derechos y sus
intenciones de involucramiento en actividades comunitarias y políticas
futuras (\citeproc{ref-isac_contribution_2013}{{Isac et~al.}, 2013};
\citeproc{ref-jerome_systematic_2024}{Jerome et~al., 2024};
\citeproc{ref-schulz_becoming_2018}{Schulz et~al., 2018}).

El género de los estudiantes muestra un comportamiento diferenciado
según la dimensión cívica analizada
(\citeproc{ref-abendschon_political_2017}{Abendschön \& Tausendpfund,
2017}; \citeproc{ref-schulz_becoming_2018}{Schulz et~al., 2018}). En las
evaluaciones internacionales, las mujeres obtienen, en promedio, mejores
resultados en conocimiento cívico que sus pares masculinos
(\citeproc{ref-schulz_becoming_2018}{Schulz et~al., 2018}). Asimismo,
las estudiantes reportan una adscripción significativamente más alta a
los principios de igualdad de derechos y una mayor predisposición a
participar en actividades comunitarias
(\citeproc{ref-abendschon_political_2017}{Abendschön \& Tausendpfund,
2017}; \citeproc{ref-resh_sense_2017}{Resh \& Sabbagh, 2017};
\citeproc{ref-schulz_becoming_2018}{Schulz et~al., 2018}). En cambio,
los hombres suelen manifestar niveles ligeramente superiores de
conocimiento o interés en áreas específicas de participación política
formal en ciertos contextos iniciales de desarrollo
(\citeproc{ref-abendschon_political_2017}{Abendschön \& Tausendpfund,
2017}).

A nivel escuela, el nivel socioeconomimco de la escuela establecimiento
influye directamente en las competencias de los estudiantes,
observándose que las escuelas que concentran poblaciones de mayor
estatus promedio registran niveles superiores de conocimiento y
expectativas de participación
(\citeproc{ref-jerome_systematic_2024}{Jerome et~al., 2024};
\citeproc{ref-schulz_becoming_2018}{Schulz et~al., 2018}). No obstante,
el contexto escolar también modera la relación entre las creencias
individuales y las actitudes normativas: en escuelas de alto rendimiento
académico promedio (SIMCE) o de mayor nivel socioeconómico, la relación
positiva entre la percepción meritocrática individual y la aceptación de
la justicia de mercado se atenúa
(\citeproc{ref-castillo_socialization_2024}{Castillo et~al., 2024}).

\subsection{Metodo}\label{metodo}

\subsubsection{Datos}\label{datos}

El analisis esta basado en la informacion proporcionada por el la
encuesta EDUMER, {[}rellenar con informacion tecnicade la ecuesta.{]}

\subsubsection{Analisis}\label{analisis}

En el presente estudio empleará dos tipos de análisis. En primer lugar,
se realizarán análisis descriptivos como gráficos univariados y
bivariados de los indicadores principales, así como también
correlaciones para analizar el comportamiento de las variables a nivel
muestral. Estos análisis entregarán un panorama previo respecto al
comportamiento de las variables y al cumplimiento de los supuestos
necesarios para los modelos posteriores.

En segundo lugar, considerando la naturaleza longitudinal y anidada de
los datos; observaciones (mediciones repetidas) anidadas en estudiantes,
y estudiantes a su vez anidados en colegios, se efectuará un análisis
multivariado empleando como técnica modelos de regresión jerárquica.
Esta estrategia de análisis resulta pertinente por dos motivos. Primero,
permite corregir la no independencia de las observaciones propia de esta
estructura de datos, la cual vulnera el supuesto de independencia de los
residuos característico de los modelos de regresión convencionales
(OLS): las mediciones repetidas de un mismo estudiante, así como las de
estudiantes que comparten colegio, tienden a estar correlacionadas entre
sí. Segundo, el enfoque multinivel permite capturar explícitamente el
efecto de variables contextuales de colegio, es decir, características
propias del establecimiento que podrían incidir en los resultados
observados. De este modo, los modelos jerárquicos permiten descomponer
la varianza total del fenómeno en los distintos niveles de anidamiento
(ocasión de medición, estudiante y colegio), estimando de forma
simultánea efectos fijos asociados a predictores de interés en cada
nivel, incluyendo los de nivel colegio y efectos aleatorios (por
definir) asociados a la variabilidad entre estudiantes y entre colegios.

Se calcularán dos conjuntos de modelos jerárquicos, uno para cada
variable dependiente del estudio, ambas construidas como índices: el
primero utilizando como variable dependiente un índice de conocimiento
cívico, y el segundo utilizando como variable dependiente un índice de
actitudes cívicas.

\bookmarksetup{startatroot}

\chapter{Antecedentes conceptuales y
empíricos}\label{antecedentes-conceptuales-y-empuxedricos}

Sintetiza la literatura relevante. Usa referencias como
(\citeproc{ref-desmond2019poverty}{\textbf{desmond2019poverty?}}).

\section{Marco conceptual}\label{marco-conceptual}

Lorem ipsum dolor sit amet, consectetur adipiscing elit. Integer at
nulla a ligula volutpat dictum. Sed dictum hendrerit augue, vitae
efficitur nisl hendrerit nec. Lorem ipsum dolor sit amet, consectetur
adipiscing elit. Integer at nulla a ligula volutpat dictum. Sed dictum
hendrerit augue, vitae efficitur nisl hendrerit nec. Lorem ipsum dolor
sit amet, consectetur adipiscing elit. Integer at nulla a ligula
volutpat dictum. Sed dictum hendrerit augue, vitae efficitur nisl
hendrerit nec. Lorem ipsum dolor sit amet, consectetur adipiscing elit.
Integer at nulla a ligula volutpat dictum. Sed dictum hendrerit augue,
vitae efficitur nisl hendrerit nec.

Lorem ipsum dolor sit amet, consectetur adipiscing elit. Integer at
nulla a ligula volutpat dictum. Sed dictum hendrerit augue, vitae
efficitur nisl hendrerit nec. Lorem ipsum dolor sit amet, consectetur
adipiscing elit. Integer at nulla a ligula volutpat dictum. Sed dictum
hendrerit augue, vitae efficitur nisl hendrerit nec. Lorem ipsum dolor
sit amet, consectetur adipiscing elit. Integer at nulla a ligula
volutpat dictum. Sed dictum hendrerit augue, vitae efficitur nisl
hendrerit nec. Lorem ipsum dolor sit amet, consectetur adipiscing elit.
Integer at nulla a ligula volutpat dictum. Sed dictum hendrerit augue,
vitae efficitur nisl hendrerit nec.

Lorem ipsum dolor sit amet, consectetur adipiscing elit. Integer at
nulla a ligula volutpat dictum. Sed dictum hendrerit augue, vitae
efficitur nisl hendrerit nec. Lorem ipsum dolor sit amet, consectetur
adipiscing elit. Integer at nulla a ligula volutpat dictum. Sed dictum
hendrerit augue, vitae efficitur nisl hendrerit nec. Lorem ipsum dolor
sit amet, consectetur adipiscing elit. Integer at nulla a ligula
volutpat dictum. Sed dictum hendrerit augue, vitae efficitur nisl
hendrerit nec. Lorem ipsum dolor sit amet, consectetur adipiscing elit.
Integer at nulla a ligula volutpat dictum. Sed dictum hendrerit augue,
vitae efficitur nisl hendrerit nec.

\subsection{Ejemplos APA 7}\label{ejemplos-apa-7}

Según (\citeproc{ref-desmond2019poverty}{\textbf{desmond2019poverty?}}),
los patrones de segregación residencial configuran los límites de
oportunidad urbana y permiten introducir la sintaxis de cita narrativa
dentro del flujo del texto.

Los matices culturales del gusto pueden reforzar desigualdades cuando se
institucionalizan en jerarquías de prestigio académico
(\citeproc{ref-bourdieu2012}{\textbf{bourdieu2012?}}), lo que ilustra la
cita parentética con número de página.

Para contrastar hallazgos se pueden reunir múltiples referencias en un
mismo paréntesis analítico
(\citeproc{ref-desmond2019poverty}{\textbf{desmond2019poverty?}};
\citeproc{ref-xie2017bookdown}{\textbf{xie2017bookdown?}}), destacando
cómo distintos diseños metodológicos convergen en un mecanismo causal
compartido.

\begin{quote}
Las aspiraciones colectivas se negocian constantemente en espacios
cotidianos que traducen capital cultural en prácticas concretas,
reescribiendo así el mapa de posibilidades percibidas por cada agente
social en contextos urbanos desiguales. Este entramado simbólico orienta
aspiraciones futuras y limita adaptaciones comunitarias.
(\citeproc{ref-bourdieu2012}{\textbf{bourdieu2012?}})
\end{quote}

\section{Resumen del capítulo}\label{resumen-del-capuxedtulo}

\begin{quote}
Este capítulo ofreció una revisión breve de los conceptos fundamentales
y estableció proposiciones para guiar el análisis empírico posterior.
\end{quote}

\bookmarksetup{startatroot}

\chapter{Metodología}\label{metodologuxeda}

\section{Datos}\label{datos-1}

Texto de muestra para describir las fuentes de datos utilizadas.

\begin{Shaded}
\begin{Highlighting}[]
\CommentTok{\# Carga paquetes de ejemplo}
\FunctionTok{library}\NormalTok{(tidyverse)}
\end{Highlighting}
\end{Shaded}

\begin{verbatim}
      mpg             cyl             disp             hp       
 Min.   :10.40   Min.   :4.000   Min.   : 71.1   Min.   : 52.0  
 1st Qu.:15.43   1st Qu.:4.000   1st Qu.:120.8   1st Qu.: 96.5  
 Median :19.20   Median :6.000   Median :196.3   Median :123.0  
 Mean   :20.09   Mean   :6.188   Mean   :230.7   Mean   :146.7  
 3rd Qu.:22.80   3rd Qu.:8.000   3rd Qu.:326.0   3rd Qu.:180.0  
 Max.   :33.90   Max.   :8.000   Max.   :472.0   Max.   :335.0  
      drat             wt             qsec             vs        
 Min.   :2.760   Min.   :1.513   Min.   :14.50   Min.   :0.0000  
 1st Qu.:3.080   1st Qu.:2.581   1st Qu.:16.89   1st Qu.:0.0000  
 Median :3.695   Median :3.325   Median :17.71   Median :0.0000  
 Mean   :3.597   Mean   :3.217   Mean   :17.85   Mean   :0.4375  
 3rd Qu.:3.920   3rd Qu.:3.610   3rd Qu.:18.90   3rd Qu.:1.0000  
 Max.   :4.930   Max.   :5.424   Max.   :22.90   Max.   :1.0000  
       am              gear            carb      
 Min.   :0.0000   Min.   :3.000   Min.   :1.000  
 1st Qu.:0.0000   1st Qu.:3.000   1st Qu.:2.000  
 Median :0.0000   Median :4.000   Median :2.000  
 Mean   :0.4062   Mean   :3.688   Mean   :2.812  
 3rd Qu.:1.0000   3rd Qu.:4.000   3rd Qu.:4.000  
 Max.   :1.0000   Max.   :5.000   Max.   :8.000  
\end{verbatim}

\section{Variables}\label{variables}

Detalle de las variables incluidas en el análisis, así como escalas de
medición.

\begin{longtable}[]{@{}
  >{\raggedright\arraybackslash}p{(\linewidth - 10\tabcolsep) * \real{0.4085}}
  >{\centering\arraybackslash}p{(\linewidth - 10\tabcolsep) * \real{0.0563}}
  >{\centering\arraybackslash}p{(\linewidth - 10\tabcolsep) * \real{0.1268}}
  >{\centering\arraybackslash}p{(\linewidth - 10\tabcolsep) * \real{0.1690}}
  >{\centering\arraybackslash}p{(\linewidth - 10\tabcolsep) * \real{0.1127}}
  >{\centering\arraybackslash}p{(\linewidth - 10\tabcolsep) * \real{0.1268}}@{}}
\caption{Estadísticas descriptivas por variable
(ejemplo)}\tabularnewline
\toprule\noalign{}
\begin{minipage}[b]{\linewidth}\raggedright
Variable
\end{minipage} & \begin{minipage}[b]{\linewidth}\centering
N
\end{minipage} & \begin{minipage}[b]{\linewidth}\centering
Media
\end{minipage} & \begin{minipage}[b]{\linewidth}\centering
Desv. Est.
\end{minipage} & \begin{minipage}[b]{\linewidth}\centering
Min
\end{minipage} & \begin{minipage}[b]{\linewidth}\centering
Max
\end{minipage} \\
\midrule\noalign{}
\endfirsthead
\toprule\noalign{}
\begin{minipage}[b]{\linewidth}\raggedright
Variable
\end{minipage} & \begin{minipage}[b]{\linewidth}\centering
N
\end{minipage} & \begin{minipage}[b]{\linewidth}\centering
Media
\end{minipage} & \begin{minipage}[b]{\linewidth}\centering
Desv. Est.
\end{minipage} & \begin{minipage}[b]{\linewidth}\centering
Min
\end{minipage} & \begin{minipage}[b]{\linewidth}\centering
Max
\end{minipage} \\
\midrule\noalign{}
\endhead
\bottomrule\noalign{}
\endlastfoot
Ingreso per cápita (mensual) & 32 & 20.091 & 6.027 & 10.400 & 33.900 \\
Capital social barrial & 32 & 3.217 & 0.978 & 1.513 & 5.424 \\
Precariedad laboral & 32 & 146.688 & 68.563 & 52.000 & 335.000 \\
Capital institucional & 32 & 3.597 & 0.535 & 2.760 & 4.930 \\
Participación comunitaria & 32 & 17.849 & 1.787 & 14.500 & 22.900 \\
\end{longtable}

\section{Métodos}\label{muxe9todos}

Breve descripción de los procedimientos aplicados para procesar y
analizar la información.

\bookmarksetup{startatroot}

\chapter{Resultados}\label{resultados}

\section{Analisis descriptivo}\label{analisis-descriptivo}

Resumen de las estadisticas descriptivas clave y visualizaciones de
apoyo.

\begin{Shaded}
\begin{Highlighting}[]
\ControlFlowTok{if}\NormalTok{ (}\FunctionTok{requireNamespace}\NormalTok{(}\StringTok{"ggplot2"}\NormalTok{, }\AttributeTok{quietly =} \ConstantTok{TRUE}\NormalTok{)) \{}
  \FunctionTok{library}\NormalTok{(ggplot2)}
  \FunctionTok{ggplot}\NormalTok{(mtcars, }\FunctionTok{aes}\NormalTok{(wt, mpg)) }\SpecialCharTok{+} \FunctionTok{geom\_point}\NormalTok{()}
\NormalTok{\}}
\end{Highlighting}
\end{Shaded}

\begin{figure}[H]

\caption{Relacion X\^{}2}

{\centering \pandocbounded{\includegraphics[keepaspectratio]{04-resultados_files/figure-pdf/grafico-ejemplo-1.pdf}}

}

\end{figure}%

\section{Modelos}\label{modelos}

Presentacion de los resultados de los modelos propuestos, con
interpretacion y conclusiones.

\begin{longtable}[]{@{}
  >{\raggedright\arraybackslash}p{(\linewidth - 12\tabcolsep) * \real{0.2738}}
  >{\centering\arraybackslash}p{(\linewidth - 12\tabcolsep) * \real{0.0952}}
  >{\centering\arraybackslash}p{(\linewidth - 12\tabcolsep) * \real{0.1310}}
  >{\centering\arraybackslash}p{(\linewidth - 12\tabcolsep) * \real{0.1071}}
  >{\centering\arraybackslash}p{(\linewidth - 12\tabcolsep) * \real{0.1071}}
  >{\centering\arraybackslash}p{(\linewidth - 12\tabcolsep) * \real{0.0714}}
  >{\centering\arraybackslash}p{(\linewidth - 12\tabcolsep) * \real{0.2143}}@{}}
\caption{Resumen comparativo de modelos}\tabularnewline
\toprule\noalign{}
\begin{minipage}[b]{\linewidth}\raggedright
Predictor
\end{minipage} & \begin{minipage}[b]{\linewidth}\centering
Coef.
\end{minipage} & \begin{minipage}[b]{\linewidth}\centering
Std. Err.
\end{minipage} & \begin{minipage}[b]{\linewidth}\centering
t-valor
\end{minipage} & \begin{minipage}[b]{\linewidth}\centering
p-valor
\end{minipage} & \begin{minipage}[b]{\linewidth}\centering
Sig.
\end{minipage} & \begin{minipage}[b]{\linewidth}\centering
IC 95 \%
\end{minipage} \\
\midrule\noalign{}
\endfirsthead
\toprule\noalign{}
\begin{minipage}[b]{\linewidth}\raggedright
Predictor
\end{minipage} & \begin{minipage}[b]{\linewidth}\centering
Coef.
\end{minipage} & \begin{minipage}[b]{\linewidth}\centering
Std. Err.
\end{minipage} & \begin{minipage}[b]{\linewidth}\centering
t-valor
\end{minipage} & \begin{minipage}[b]{\linewidth}\centering
p-valor
\end{minipage} & \begin{minipage}[b]{\linewidth}\centering
Sig.
\end{minipage} & \begin{minipage}[b]{\linewidth}\centering
IC 95 \%
\end{minipage} \\
\midrule\noalign{}
\endhead
\bottomrule\noalign{}
\endlastfoot
Constante & 37.227 & 1.599 & 23.28 & \textless0.001 & *** & {[}33.957,
40.497{]} \\
Capital social barrial & -3.878 & 0.633 & -6.13 & \textless0.001 & *** &
{[}-5.172, -2.584{]} \\
Precariedad laboral & -0.032 & 0.009 & -3.52 & 0.001 & ** & {[}-0.050,
-0.013{]} \\
N & 32 & & & & & \\
R\^{}2 & 0.827 & & & & & \\
R\^{}2 ajustado & 0.815 & & & & & \\
\end{longtable}

\bookmarksetup{startatroot}

\chapter{Discusión}\label{discusiuxf3n}

Interpreta resultados, limitaciones y aportes.

\bookmarksetup{startatroot}

\chapter{Conclusiones}\label{conclusiones}

Síntesis y proyecciones futuras.

\bookmarksetup{startatroot}

\chapter{Referencias}\label{referencias}

\protect\phantomsection\label{refs}
\begin{CSLReferences}{1}{0}
\bibitem[\citeproctext]{ref-abendschon_political_2017}
Abendschön, S., \& Tausendpfund, M. (2017). Political {Knowledge} of
{Children} and the {Role} of {Sociostructural Factors}. \emph{American
Behavioral Scientist}, \emph{61}(2), 204-221.
\url{https://doi.org/10.1177/0002764216689122}

\bibitem[\citeproctext]{ref-batruch_belief_2023}
Batruch, A., Jetten, J., Van De Werfhorst, H., Darnon, C., \& Butera, F.
(2023). Belief in {School Meritocracy} and the {Legitimization} of
{Social} and {Income Inequality}. \emph{Social Psychological and
Personality Science}, \emph{14}(5), 621-635.
\url{https://doi.org/10.1177/19485506221111017}

\bibitem[\citeproctext]{ref-campbell_voice_2008}
Campbell, D. E. (2008). Voice in the {Classroom}: {How} an {Open
Classroom Climate Fosters Political Engagement Among Adolescents}.
\emph{Political Behavior}, \emph{30}(4), 437-454.
\url{https://doi.org/10.1007/s11109-008-9063-z}

\bibitem[\citeproctext]{ref-canalesceron_sujeto_2022}
Canales Cerón, M., Orellana Calderón, V. S., Guajardo Mañán, F., Canales
Cerón, M., Orellana Calderón, V. S., \& Guajardo Mañán, F. (2022).
{Sujeto y cotidiano en la era neoliberal: el caso de la educación
chilena}. \emph{Revista mexicana de ciencias políticas y sociales},
\emph{67}(244), 285-307.
\url{https://doi.org/10.22201/fcpys.2448492xe.2022.244.70386}

\bibitem[\citeproctext]{ref-castillo_socialization_2024}
Castillo, J. C., Salgado, M., Carrasco, K., \& Laffert, A. (2024). The
{Socialization} of {Meritocracy} and {Market Justice Preferences} at
{School}. \emph{Societies}, \emph{14}(11), 214.
\url{https://doi.org/10.3390/soc14110214}

\bibitem[\citeproctext]{ref-castillo_meritocracia_2019}
Castillo, J. C., Torres, A., Atria, J., \& Maldonado, L. (2019).
Meritocracia y Desigualdad Económica: {Percepciones}, Preferencias e
Implicancias. \emph{Revista Internacional de Sociología}, \emph{77}(1),
e117. \url{https://doi.org/10.3989/ris.2019.77.1.17.114}

\bibitem[\citeproctext]{ref-castillo_social_2014}
Castillo, J., Miranda, D., Bonhomme, M., Cox, C., \& Bascopé, M. (2014).
Social Inequality and Changes in Students' Expected Political
Participation in {Chile}. \emph{Education, Citizenship and Social
Justice}, \emph{9}, 140-156.
\url{https://doi.org/10.1177/1746197914520650}

\bibitem[\citeproctext]{ref-crawford_schooling_2010}
Crawford, K. (2010). Schooling, {Citizenship} and the {Myth} of the
{Meritocracy}. \emph{Citizenship, Social and Economics Education},
\emph{9}(1), 3-13. \url{https://doi.org/10.2304/csee.2010.9.1.3}

\bibitem[\citeproctext]{ref-darnon_where_2018}
Darnon, C., Wiederkehr, V., Dompnier, B., \& Martinot, D. (2018).
{«{Where} There Is a Will, There Is a Way»}: {Belief} in School
Meritocracy and the Social-class Achievement Gap. \emph{British Journal
of Social Psychology}, \emph{57}(1), 250-262.
\url{https://doi.org/10.1111/bjso.12214}

\bibitem[\citeproctext]{ref-ferre_welfare_2023}
Ferre, J. C. (2023). Welfare Regimes in Twenty-First-Century {Latin
America}. \emph{Journal of International and Comparative Social Policy},
\emph{39}(2), 101-127. \url{https://doi.org/10.1017/ics.2023.16}

\bibitem[\citeproctext]{ref-flanagan_adolescents_2014}
Flanagan, C. A., Kim, T., Pykett, A., Finlay, A., Gallay, E. E., \&
Pancer, M. (2014). Adolescents' Theories about Economic Inequality:
{Why} Are Some People Poor While Others Are Rich? \emph{Developmental
Psychology}, \emph{50}(11), 2512-2525.
\url{https://doi.org/10.1037/a0037934}

\bibitem[\citeproctext]{ref-isac_contribution_2013}
{Isac, M. M., Maslowski, R., Creemers, B., \& van der Werf, G.} (2013).
The Contribution of Schooling to Secondary-School Students' Citizenship
Outcomes across Countries. \emph{School Effectiveness and School
Improvement}. \url{https://doi.org/10.1080/09243453.2012.751035}

\bibitem[\citeproctext]{ref-jerome_systematic_2024}
Jerome, L., Hyder, F., Hilal, Y., \& Kisby, B. (2024). A Systematic
Literature Review of Research Examining the Impact of Citizenship
Education on Active Citizenship Outcomes. \emph{Review of Education},
\emph{12}(2), e3472. \url{https://doi.org/10.1002/rev3.3472}

\bibitem[\citeproctext]{ref-lee_processes_2012}
Lee, N.-J., Shah, D. V., \& McLeod, J. M. (2012). Processes of
{Political Socialization}. \emph{Communication Research}.
\url{https://doi.org/10.1177/0093650212436712}

\bibitem[\citeproctext]{ref-luismaldonado_meritocracia_2021}
Maldonado, L., Castillo, J. C., Julio Iturra-Sanhueza, Rivas, F. J. M.,
\& Atria, J. (2021). \emph{Meritocracia y Redistribución En {Chile}:
Señales de La Opinión Pública}.
\url{https://doi.org/10.17605/OSF.IO/G4EK8}

\bibitem[\citeproctext]{ref-palet_desiguales_2017}
Palet, A., Aguirre, P. de, \& PNUD Chile (Eds.). (2017).
\emph{{Desiguales: orígenes, cambios y desafíos de la brecha social en
Chile}}. PNUD : Uqbar Editores.

\bibitem[\citeproctext]{ref-_pdf_c}
({PDF}) {Desigualdad} y {Conocimiento Cívico}: {Chile} En {Comparación
Internacional}. (s.~f.). En \emph{{ResearchGate}}.

\bibitem[\citeproctext]{ref-resh_sense_2010}
Resh, N. (2010). Sense of Justice about Grades in School: Is It
Stratified like Academic Achievement? \emph{Social Psychology of
Education}, \emph{13}(3), 313-329.
\url{https://doi.org/10.1007/s11218-010-9117-z}

\bibitem[\citeproctext]{ref-resh_sense_2017}
Resh, N., \& Sabbagh, C. (2017). Sense of Justice in School and Civic
Behavior. \emph{Social Psychology of Education}, \emph{20}(2), 387-409.
\url{https://doi.org/10.1007/s11218-017-9375-0}

\bibitem[\citeproctext]{ref-schulz_becoming_2018}
Schulz, W., Ainley, J., Fraillon, J., Losito, B., Agrusti, G., \&
Friedman, T. (2018). \emph{Becoming {Citizens} in a {Changing World}:
{IEA International Civic} and {Citizenship Education Study} 2016
{International Report}}. Springer International Publishing.
\url{https://doi.org/10.1007/978-3-319-73963-2}

\bibitem[\citeproctext]{ref-somma_no_2020}
Somma, N., Bargsted, M., Disi Pavlic, R., \& Medel, R. (2020). No Water
in the Oasis: The {Chilean Spring} of 2019--2020. \emph{Social Movement
Studies}, \emph{20}, 1-8.
\url{https://doi.org/10.1080/14742837.2020.1727737}

\bibitem[\citeproctext]{ref-verger_recontextualizacion_2016}
Verger, A., Bonal, X., \& Zancajo, A. (2016). Recontextualización de
Políticas y (Cuasi)Mercados Educativos. {Un} Análisis de Las Dinámicas
de Demanda y Oferta Escolar En {Chile}. \emph{Education Policy Analysis
Archives}, \emph{24}, 27. \url{https://doi.org/10.14507/epaa.24.2098}

\bibitem[\citeproctext]{ref-wiederkehr_belief_2015}
Wiederkehr, V., Bonnot, V., Krauth-Gruber, S., \& Darnon, C. (2015).
Belief in School Meritocracy as a System-Justifying Tool for Low Status
Students. \emph{Frontiers in Psychology}, \emph{6}.
\url{https://doi.org/10.3389/fpsyg.2015.01053}

\bibitem[\citeproctext]{ref-young_triunfo_}
Young, M. (s.~f.). \emph{{Triunfo de la meritocracia}}.

\end{CSLReferences}

\cleardoublepage
\phantomsection
\addcontentsline{toc}{part}{Apéndices}
\appendix

\chapter{Instrumentos de
levantamiento}\label{instrumentos-de-levantamiento}

Incluye cuestionarios, guias de entrevista, etc.

\chapter{Tablas adicionales}\label{tablas-adicionales}

Tablas extendidas, tests de robustez, etc.


\backmatter

% --- Cuerpo del libro (capítulos) en arábigos ---
\mainmatter
\pagestyle{scrheadings}

% --- (Opcional) Estilo del índice general ---
% \addtocontents{toc}{\protect\thispagestyle{plain}}


\end{document}
